% Layout fixture: geometry and mechanics notes whose signs TeX draws from pieces: arrows
% over points (\overrightarrow, a shaft of minus glyphs glued to a head), a left arrow
% next to a minus, \cong (a tilde over an equals sign) and \models (a bar and an equals
% sign), and vector products written with \vec. Used by tests/test_sign_check.py.
% Written for ReadableSTEM's tests; no third-party text.
\documentclass[11pt]{article}
% Computer Modern text fonts (CMR, CMBX), as pdfLaTeX writes them; XeTeX defaults to Latin Modern.
\usepackage[OT1]{fontenc}
\renewcommand{\rmdefault}{cmr}
\usepackage{amsmath,amssymb}
\usepackage[margin=1in]{geometry}
\pagestyle{empty}

\begin{document}

\section*{Vectors between points}

The displacement from $A$ to $B$ is the difference of the two position vectors
measured from the origin,
\[
  \overrightarrow{AB} = \overrightarrow{OB} - \overrightarrow{OA}.
\]
Reversing the order of the points reverses the vector:
\[
  \overrightarrow{AB} = -\overrightarrow{BA}.
\]
Two triangles with equal sides are congruent, which we write
\[
  \triangle ABC \cong \triangle DEF.
\]

\section*{Updates and logic}

A loop that counts down starts with the assignment
\[
  x \gets -1
\]
and a premise set that entails a formula is written
\[
  \Gamma \models \varphi.
\]

\section*{Work and torque}

The work done by a constant force over a displacement is a dot product,
\[
  W = \vec{F} \cdot \vec{d},
\]
while the torque about the origin is a cross product,
\[
  \vec{\tau} = \vec{r} \times \vec{F}.
\]

\end{document}
